% Two routes to the same row: node(id:), built in chapter 4, and the
% _entities(representations:) reference resolver this chapter adds.
%
% Same idiom as figures/tikz/ch01-one-field.tex and
% figures/tikz/ch05-the-join-moves.tex: each lane is a timeline read
% left to right, ticks mark stages, and a brace under the stage that
% matters carries a one-line label. The four ticks sit at the same x
% in both lanes, because the argument is that the two routes do the
% same work at the same points and differ only in what they are
% handed and in where the decoding happens.
%
% The one asymmetry the figure draws: in lane A, decoding is its own
% stage, tick 2, run by Hot Chocolate before any resolver is entered.
% Lane B's tick 2 is the same point in the sequence and nothing happens
% there, which its label says; the id string arrives whole and the
% reference resolver, tick 3, decodes it itself. The brace therefore
% lands under tick 2 in lane A and tick 3 in lane B, and that shift is
% what the eye should catch first.
%
% Lane B's SELECT elides the middle four columns and states the count
% in words, keeping the first and last column named. Both lanes cost
% one statement for one key; neither is drawn as cheaper.
%
% No quotation marks anywhere in this file. SPEC decision 30 puts the
% book in strict Ascii mode and the prose gate reads a literal " as a
% quote character, while \" is LaTeX's diaeresis accent and fails to
% compile inside a node. The labels are written so that neither is
% needed: identifiers appear bare.
%
% Bare tikzpicture (house style): the figure environment, caption and
% label live at the call site in the section file.
\begin{tikzpicture}[
  x=1mm, y=1mm,
  every node/.append style={font=\sffamily\scriptsize},
  axis/.style={draw, line width=0.4pt,
    -{Straight Barb[length=1.4mm, width=1.6mm]}},
  tick/.style={draw, line width=0.4pt},
  event/.style={anchor=south, align=center, text width=24mm,
    inner sep=1pt},
  span/.style={draw, line width=0.4pt},
  spanlabel/.style={anchor=north, align=center, inner sep=2pt},
  lane/.style={anchor=west, font=\sffamily\scriptsize\itshape},
]

% ------------------------------------------------------------- lane A
\node[lane] at (0,20)
  {A. the node route, asking for one field};

\draw[axis] (0,0) -- (150,0);
\foreach \x in {16,58,100,138}{
  \draw[tick] (\x,-1.5) -- (\x,1.5);
}

\node[event, text width=30mm] at (16,3)
  {client sends\\node(id: U2Vzc2lvbjox),\\selecting title};
\node[event, text width=28mm] at (58,3)
  {the id is decoded\\before the resolver\\is entered};
\node[event, text width=30mm] at (100,3)
  {node resolver is given\\the integer 1 and a\\QueryContext};
\node[event, text width=26mm] at (138,3)
  {SELECT s.Title:\\one column};

\draw[span] (44,-4) -- (44,-7) -- (72,-7) -- (72,-4);
\node[spanlabel] at (58,-8)
  {decoded here, by the framework,\\as a stage of its own};

% ------------------------------------------------------------- lane B
\begin{scope}[shift={(0,-46)}]
  \node[lane] at (0,20)
    {B. the entity route, asking for the same field};

  \draw[axis] (0,0) -- (150,0);
  \foreach \x in {16,58,100,138}{
    \draw[tick] (\x,-1.5) -- (\x,1.5);
  }

  \node[event, text width=32mm] at (16,3)
    {router sends one\\representation:\\\_\_typename Session,\\id
     U2Vzc2lvbjox};
  \node[event, text width=26mm] at (58,3)
    {nothing decodes\\anything};
  \node[event, text width=30mm] at (100,3)
    {reference resolver is\\given the string, and\\no QueryContext};
  \node[event, text width=26mm] at (138,3)
    {SELECT s.Id, ...,\\s.Title:\\six columns};

  \draw[span] (86,-4) -- (86,-7) -- (114,-7) -- (114,-4);
  \node[spanlabel] at (100,-8)
    {decoded here, by hand, inside\\the resolver you wrote};
\end{scope}

\end{tikzpicture}
